
\section{Introduction}
 
In the past two centuries, over half of the world's nations have transitioned to democracy.  The predominant view in political economy and political science argues these democratizations stem from intrinsic conflicts among different political or social classes and the owners of the means of production \citep*{Marx1848,Lipset1959,Moore1966, Boix2003,Acemoglu2006,Ansell2010}.  Consensus on this point, however, remains elusive, stemming in part from two broad empirical challenges.  First, there is little evidence on whether, {\it ex ante}, elites consider redistribution to be the central risk they face in democratization. Second, conditional on redistribution risk being central,  there is little understanding of which among the various forms of political and economic redistribution that democracy might entail are most important \citep{Acemoglu2015}.

This paper tests whether redistribution is central to democratizations in a way that addresses both of these issues: examining stock market prices during democratizations. Since asset prices disproportionately reflect the expectations and preferences of mostly wealthy capital holders---especially in autocratic countries---they are an ideal source for understanding the risk the elites perceive from the democratization process in real time. 

How do financial markets respond when democratization becomes more likely? Using a panel of equity data that covers 90 countries over 200 years, I show that stock market valuations fall substantially when transitions to democracy are more likely.  In the data,  I document that this decline is similar in magnitude to what we observe in financial crises, suggesting that these periods are associated with increased systematic risk to investors.

To understand whether the risk of redistribution drives this result, however, two key empirical challenges must be addressed. First,  it is essential to tackle potential endogeneity concerns by ruling out other common factors that could simultaneously affect democratizations and financial markets, and provide evidence that the ancillary effects of democratization---for example, political instability or violence---are not driving the results. Second, it is necessary to show that the primary driver of the asset pricing response is redistribution risk. This requires showing that redistribution indeed follows successful democratizations,  and that it is substantial enough to rationalize the observed market responses.  The remainder of the paper provides evidence of these two central points.

The first part of the paper uses two main strategies to document that democratizations indeed drive the negative stock market response.  First, I directly show that several potential first-order channels are unlikely to be driving the results.  For example, democratizations could coincide with an increase in macroeconomic risk which would tend to drive down stock valuations. However, this is not borne out in the data. GDP or aggregate dividend growth do not fall in the 5 years after a democratization starts, nor do the distributions of GDP or aggregate consumption growth change.  A rise in generic political risk cannot fully explain the results either. Other periods of high political risk like international political crises, autocratizations---transitions from democracy to autocracy---and other regime changes exhibit substantially smaller stock market responses when compared to democratizations.\footnote{A lengthy online appendix provides evidence against several other potential explanations, like increased violence, the increased probability of adverse financial or macroeconomic events, increased revolution risk, large capital outflows, and general uncertainty shocks.}

Second, I show democratizations increase risk premia using two identification strategies.  The first strategy uses exogenous variation in the probability of a successful democratization emanating from a shift in Catholic church doctrine in favor of democracy from 1959 to 1963. This shift particularly impacted majority Catholic autocracies. \cite{Huntington1991} labels the shift as one of the main reasons the third wave of democratization of the 1970s, 1980s, and early 1990s occurred and why it was concentrated in majority Catholic autocracies. Consistent with this narrative,  I show that indices denoting the threat to the governing regime posed by civil society organizations and the size and frequency of democratic protests rose dramatically in majority Catholic autocracies compared to non-Catholic autocracies. This indicates that the doctrinal shift materially changed political realities on the ground in majority Catholic autocracies. 

Using a difference-in-differences approach, this quasi-natural experiment is associated with a \tabSevenLow\ to \tabSevenHigh\ percentage point increase in average excess stock returns for majority Catholic autocracies depending on the specification. The results display no pre-trends and are robust to various sample windows, the exclusion of outliers, and different estimation techniques. They also cement the link between an increase in risk premia and an increase in the probability of a successful democratic transition.

The second identification strategy follows \cite*{Acemoglu2019} and uses regional waves of democratization as exogenous variation in the likelihood of a successful democratization. Democratic progress from regional waves is more likely to be driven by external pressure and is, therefore, exogenous to the local macroeconomic conditions of a particular country.  Democratic movements coming from regional waves also see a substantial increase in dividend yields in both a reduced form and instrumental variables approach.

The second part of the paper investigates whether a rise in redistribution risk can explain the negative stock market reaction to democratizations.  Comparing successful and failed democratizations, I find that democratization redistributes resources in two ways. First, it increases explicit redistribution by raising the size of the public sector and lowering income inequality. On average, government revenue-GDP ratios rise by \tabEightOneCE\ percentage points, Gini coefficients decline by \tabEightThreeCE\ percentage points, and the labor share of GDP rises by \tabEightFourCE\ percentage points in the 20 years after a successful democratization. Second, successful democratizations also increase tacit redistribution. For example,  autocracies allocate a greater share of government spending to elites \citep{Tullock1986}. They also provide more protection to incumbent firms from new entrants \citep{Perotti2006,Bravo2021}. I find that, during successful democratizations, bribery and corruption indices fall while pro-competitive regulation and net entry of new firms rise. Since this also redistributes resources away from autocratic elites, it could also play an important role in the asset pricing results.

To understand whether the redistribution in the data is quantitatively large enough to explain the asset pricing results, I calibrate a model of democratic transitions in the style of \cite{Acemoglu2006} embedded within a standard asset pricing framework. Like in  \citeauthor{Acemoglu2006}, the economy starts in autocracy where the elites have all the political power, and try to avoid redistributing their income to the more numerous poor citizens. The citizens influence the policies of the elites by threatening to revolt. Revolution is costly: all the elites are killed and a fraction of resources are destroyed, making it undesirable for both sides.  This cost the citizens bear from revolution---which determines the revolutionary threat the elites face---varies over time.  If the fraction of resources destroyed is low enough, though, the citizens may prefer the revolution to autocracy.  When this happens, the elites would like to promise future redistribution.  But they cannot credibly commit to future transfers where there is little or no revolutionary threat. Here, only conceding democracy can keep the revolution off the equilibrium path, as democracy acts as a mechanism for the elites to credibly commit to future redistribution. While democracy is a much better state for the elites than the revolution, the redistribution it brings is costly, making it, nonetheless, a deleterious state for them. 

To make the model relevant to study asset prices, I add four main ingredients. First, I allow for incomplete financial markets, meaning that the elites can trade with one another in financial markets but not with the citizens.  Second, to achieve realistic asset pricing dynamics, I allow for preferences in the style of \cite{Epstein1989}. Third, I allow for multiple potential forms of redistribution that align with what we see in the data, namely, reduced inequality, increased taxes, reduced ability for the elites to skim rents from government spending, and increased economic competition. I also allow for the redistribution elites face in democracy to be uncertain.

Fourth, I modify the cost of revolution process to allow for three states: autocracy, democratization, and democracy. The new state, democratization, is one where a permanent transition to democracy becomes more likely. Since the elites price assets, uncertainty over whether a democratization will succeed---ushering in democracy and redistribution---or fail---keeping society in autocracy---increases the risk to the elites' future consumption, causing risk premia to rise. In this way, the consolidation of democracy and the redistribution of income and political power it brings, acts as a ``rare disaster'' for the elites, explaining the increased risk premia observed during democratizations in the data \citep*{Rietz1988, Barro2006a,Gabaix2012,Wachter2013}.  When calibrated to reasonable preference parameters and the redistribution observed in the data, the model explains nearly all the rise in dividend yields observed during democratizations.

The model also allows me to understand which forms of redistribution have the largest effect on asset prices. The predominant effect comes from increased economic competition and displacement risk for incumbent firms post-democratization \citep*{Garleanu2012}. This channel drives \tabTwelveComp\% of the rise in dividend yields, providing support to a theoretical literature that argues increased creative destruction and structural transformation are the primary driving forces behind higher growth after successful democratizations \citep*{Aghion2008,Aghion2014,Acemoglu2015,Bravo2021}. The remaining \tabTwelveNonComp\% of the rise in dividend yields comes from the more traditional channels of higher taxes and reduced inequality and corruption. 

A redistribution-based framework also explains the negligible stock market effect observed in autocratizations. To do this, I modify the model and allow for democracy to be reversible provided the elites are willing to risk a transition. If they succeed, society becomes an autocracy, but if they fail, they face a permanent loss of a fraction of their consumption. The key insight is that while democratization is a risk imposed on the elites, autocratization is a risk they take. Because who decides to transition differs in each case, there is an asymmetric effect on asset prices.

The elites optimally choose when to attempt autocratization, so it always improves the expected present value of their consumption. However, levered claims to this consumption---for example, the dividend claim---can still be adversely affected. In the model, dividend yields still rise because the increased risk in the event of a failed autocratization matters more than the higher payoff upon success to a risk averse investor. This also leads autocratizations with higher potential payoffs to come with larger rises in dividend yields, as the elites accept a higher penalty in the event of failure to achieve autocracy.

Taken together, these results provide powerful support for redistribution-based models of democratization. When modified to incorporate asset prices, the predictions the model generates enjoy resounding support in the data.  This helps to clear a significant hurdle in this literature. While most studies have focused on whether more democratic institutions lead to redistribution, few have substantiated whether this redistribution is large enough to constitute a major friction to democratic transitions \citep{Boix2003,Hinnerich2014,Acemoglu2015}. Better understanding this is important for the many countries still living under autocratic political institutions. It is also relevant for countries with backsliding democratic institutions, the number of which some scholars allege have increased over the last decade \citep{Diamond2015}. Insofar as reductions in democratic norms are accompanied by lower taxes, higher inequality, lower labor bargaining power, and decreased economic competition, this paper provides a model through which future autocratic movements can be interpreted.

\label{par:limitation_acknowledge_beg}
Of course, the analysis is not without limitations. It is worth highlighting two caveats. First, there is substantial heterogeneity around the average effects I document. For example, some democratizations may come with a reduction in risk premia as securing property rights dominates other channels. Second, the results above are strongest for countries transitioning to democracy with an active stock market.  This means that countries where property rights were not secure enough to foster open financial markets are scoped out of the analysis when examining valuations. One example of this is transitions from left-wing authoritarian states, like the disintegration of the Soviet Union. Since regimes of this style were not fond of capital markets, their data on asset prices do not generally exist. My analysis, therefore, helps us to better understand transitions from relatively more right-wing autocracies, where the threat of redistribution likely played a greater role.

\label{par:low_redist_evidence1}
The evidence from equity markets, however, has enough coverage to give a sense of the extent of these limitations. To better understand the first caveat, I show results for various cuts of the sample to examine where redistribution risk is the operative channel versus securing property rights. In particular, I split democratizations into high and low redistribution risk categories using information on who the most powerful political group is prior to the episode start. This exercise shows that almost all the rise in dividend yields comes from high redistribution risk democratizations.  Indeed, democratizations with low redistribution risk come with rises in dividend yields quite similar to other regime transition episodes. This is in line with the idea that not all democratizations are necessarily characterized by redistribution risk.

\label{par:low_redist_evidence2}
Related evidence in favor of this point is that nearly all the rise in dividend yields leading into democratizations is concentrated in the period after World War I, the beginning of the First Wave of Democratization.  This is in line with a narrative in the comparative politics literature about the nature of democratization before and after The Great War \citep{Luebbert1991}. Before the war democratizations were mainly agreements between the aristocracy and the burgeoning middle class, shutting out the then nascent labor movements. As such, they benefited this new capital-owning class by protecting their property rights and the \textit{status quo} between capital and labor. It is also consistent with the literature discussing the case of Britain after the Glorious Revolution documented in both \cite{North1989} and \cite{Acemoglu2006}.\footnote{For example, prior to the reforms of William III after the Glorious Revolution, property rights were not secure enough in England to allow for the formation of stock markets \citep{Brodhurst1897}. From \cite{Brodhurst1897}: ``In the time of James I, the excitements of the Stock Exchange, and the allurements of stock-brokers had not yet begun to trouble the English people. A national debt, the creation of the Venetians, was as yet unknown in England. Loans, indeed, to satisfy the necessities of State had been raised by Henry VIII and many others of the English Sovereigns; but as they never thought of repaying money which they had borrowed, and as those who were forced to lend, probably had not any expectation of seeing their property again, there was little opportunity for speculation. It was left to William III to introduce the principle that it is the duty of a State to keep faith with its creditors, and thereby to open the door to those commercial movements which were ultimately to result in the creation of the Stock Exchange.''} After World War I, however, democratizations became more labor driven, focusing on increasing labor bargaining power and reducing inequalities. Transitioning to democracy thus became more costly for the capital-owning elites, bringing higher risk premia in the transition period.

\label{par:ik_intro_evidence}
There also is evidence to suggest that the results from stock markets are applicable to most transitions, at least in the sample after World War II. Investment-capital ratios fall substantially in the five years leading into a democratization, suggesting that either risk premia are rising, or expected cashflow growth is falling, for all capital assets, depressing investment even in countries without stock markets. Since data on investment-capital ratios are present for the vast majority of democratizations that occurred after 1960, this provides strong evidence that in the recent sample democratizations do not seem to be viewed favorably by capital investors. This also shows that the results from financial markets are not necessarily limited to public equity markets.
\label{par:limitation_acknowledge_end}

Finally, while declining stock valuations following democratization might prompt concerns, it should not be interpreted as a shortcoming of democracy. The analysis above suggests the opposite: the vast majority of citizens experience notable welfare gains from democratic transitions. Instead, it hints that for markets to truly reflect the outlook of the broader macroeconomy, economic representation is paramount. The findings instead speak to a rift between Wall Street and Main Street when the goals of the wealthy and middle class come into conflict.

\paragraph{Related Literature}

This paper advances both the political economy literature around democratizations and asset pricing literature focused on rare events and political and policy risk.

My primary contributions to the political economy and democratization literatures are twofold. The first is theoretical: By adding asset prices to the seminal model in \cite*{Acemoglu2006}, this paper shows that falling asset valuations are consistent with increases in the redistribution risk faced by autocratic elites during periods of democratization. This provides a testable prediction for redistribution-based models \citep*{Boix2003,Acemoglu2006}. Moreover, the model can also assess whether the redistribution observed in the data is quantitatively large enough to explain the rise in premia. This helps clear a significant hurdle in this literature: whether the redistribution faced by the wealthy in autocracy is large enough to constitute a substantial friction to democratic transitions. 

The second is empirical. The paper provides the first evidence of the effects of democratizations on equity markets. Prior research examining the asset pricing impact of democratizations has focused on the impact on sovereign debt yields in the pre-World War I sample. Consistent with my results, it has found that suffrage extensions increase sovereign loan yields \citep{Dasgupta2021,Tuncer2022}. \cite*{Delis2019} also study the response of corporate loan spreads to democratic institutions from 1984--2014 and find that more democratic institutions are accompanied by reduced loan spreads for companies. These positive effects after transitions are not inconsistent with increased risk during the transition period, which this paper documents. Prior work has also examined the returns to politically-connected firms during regime changes. \cite{Fisman2001} finds strong negative returns for politically connected firms in Indonesia as a result of the fall of the Suharto regime. Similarly, \cite*{Acemoglu2018} find that more intense protests in Egypt after the fall of the Mubarak regime relate to lower stock market valuations for firms connected to the group currently in power. \cite*{Dube2011} find that US companies that stood to benefit from US-backed coups see high returns after the coup. My paper builds on this body of research by providing the longest time series and widest panel of equity data used to date to study the stock market impact of democratizations.

In addition to new empirical evidence on asset prices, the paper also provides a novel exercise to quantify the amount of redistribution after successful democratizations by comparing them to failed democratizations. As such, the paper compares two groups of countries that underwent a similar period of political change, but where one group experiences a sustained change and the other does not. These results, therefore, add to those reported in \cite*{Rodrik1999}, \cite*{Acemoglu2015}, and \cite*{Drautzburg2022} who measure the impact of democracy on wages, the size of the public sector, and the labor share of income. This also relates to a string of papers that study redistribution, the provision of public goods, and government spending stemming from enfranchisement episodes. These include papers studying the enfranchisement or disenfranchisement of Black Americans \citep{Husted1997,Naidu2012,Cascio2013} and women \citep{Miller2008} and various enfranchisement episodes in Western Europe \citep{Aidt2009,Hinnerich2014} and those stemming from more effective voting technology \citep{Fujiwara2015}. 

The primary contribution to the asset pricing literature is showing that large, redistributive political shocks like democratizations can act similarly to ``rare disasters'' both empirically and theoretically. In disaster models, investors are exposed to large negative shocks that manifest with some small, usually time-varying, probability \citep*{Rietz1988,Barro2006a,Gabaix2012,Wachter2013}. Investors demand compensation for holding assets exposed to these disasters, allowing these models to match key asset pricing moments. My paper adds to this literature by noting that large political risks like democratizations can come with---from the perspective of wealthy market participants---left-skewed distributional shocks which also drive asset prices.

An alternative view is offered by models where aggregate shocks affect investors differently, often through their uninsurable labor income or human capital \citep{Mankiw1986, Constantinides1996,Constantinides2017, Schmidt2016,Paron2021}.  This leads these investors to demand compensation for holding stocks allowing these models to match the level, volatility, and cross-section of asset prices. However, to generate quantitatively important asset pricing effects, these shocks need to most strongly affect the wealthy capital holders \citep*{Catherine2021}. This is the case during democratizations, as the shocks to inequality, tax policy, or political connections they bring mainly affect the wealthy.

This paper also builds on a literature examining the role of political and policy risk in asset pricing by noting that democratizations are accompanied by large increases in risk premia. \cite*{Pastor2012,Pastor2013} propose a model in which government policy uncertainty drives variation in the risk premium.  \cite*{Pastor2016} model the effect of redistributive taxation on inequality jointly with the effect on aggregate productivity and asset prices.  \cite{Pastor2021} examines how rising consumption inequality can influence to move toward populism even in a strong economy in a model in which agents are inequality averse.  Related to these papers is a literature studying the role of fluctuations in factor prices for equity prices and investment. In this context, \cite{Danthine2002} find that empirical fluctuations in the labor share combined with operating leverage can explain the unconditional level of the equity premium. \cite{Santos2005} complement this by showing that variation in the labor income to consumption ratio generated substantial time series predictability.  My paper builds on these papers by studying redistribution shocks explicitly in the context of democratizations and studying their quantitative impact on asset prices.

Empirical research on policy shocks and uncertainty has focused mostly on quantifying the effects of policy shocks in developed democracies. For example,  \cite*{Baker2016} develop an index of economic policy uncertainty and find that increases in this index are associated with greater stock price volatility and reduced investment and employment.  \cite*{Kelly2016} provide empirical support that political uncertainty is priced in the equity options market.  \cite*{Manela2017} show that variation in a text-based measure of macroeconomic and policy uncertainty co-moves with risk premia, lending credence to rare disasters theories.  Their measure of policy uncertainty also predicts future tax changes in the United States.
My paper differs from these by studying uncertainty over political institutions rather than over particular policy decisions. As such my work complements this body of research, showing that uncertainty over the institutions is also priced in financial markets.